\section{Math \& Physics Notation}
The package includes \texttt{physics} and \texttt{siunitx} for standard notation.

\subsection{Governing Equations}
The incompressible Navier-Stokes momentum equation can be written compactly using the \texttt{physics} package commands like \texttt{\textbackslash pdv} and \texttt{\textbackslash vb}:
\begin{equation} \label{eq:ns}
    \pdv{\vb{u}}{t} + (\vb{u} \cdot \grad) \vb{u} = -\frac{1}{\rho} \grad p + \nu \grad^2 \vb{u} + \vb{g}
\end{equation}

\subsection{Units}
Use \texttt{siunitx} for consistent formatting.
\begin{itemize}
    \item \textbf{Inline Quantity:} The inlet velocity is \qty{15.3}{\meter\per\second}.
    \item \textbf{Scientific Notation:} The viscosity is \qty{1.00e-6}{\meter\squared\per\second}.
    \item \textbf{Just Units:} The units of pressure are \unit{\pascal}.
\end{itemize}

\section{Tables (Tabularray)}
We use \texttt{tabularray} for professional tables. It separates content from styling.

\begin{table}[H]
    \centering
    \caption{Simulation Parameters}
    \label{tab:params}
    \begin{tblr}{
        colspec = {l c c},
        hlines,
        vlines,
        row{1} = {font=\bfseries, bg=gray9}
    }
        Parameter & Symbol & Value \\
        Density   & $\rho$ & \qty{998}{\kilo\gram\per\meter\cubed} \\
        Viscosity & $\mu$  & \qty{1.002e-3}{\pascal\second} \\
        Diameter  & $D$    & \qty{0.05}{\meter} \\
    \end{tblr}
\end{table}

\section{Code Highlighting (Minted)}
You can include Python code directly.

\subsection{Embedded Code}
Using the \texttt{minted} environment:

\begin{minted}{python}
import numpy as np

def calculate_reynolds(rho, u, d, mu):
    """
    Calculate Re = rho * u * d / mu
    """
    re = (rho * u * d) / mu
    return re

re_val = calculate_reynolds(998, 15.3, 0.05, 1.002e-3)
print(f"Reynolds Number: {re_val:.2f}")
\end{minted}

\subsection{Importing Files}
If you have an external script (e.g., \texttt{solver.py}), use this command:
\begin{verbatim}
    \inputminted{python}{solver.py}
\end{verbatim}

For scripts that span multiple pages, use the \texttt{longlisting} environment to enable page breaks while maintaining a caption:
\begin{verbatim}
    \begin{longlisting}
        \caption{Solver Implementation}
        \label{lst:solver}
        \inputminted{python}{solver.py}
    \end{longlisting}
\end{verbatim}

\section{Figures}
Standard \texttt{graphicx} support is included. References are handled by \texttt{cleveref} (e.g., see \Cref{fig:example} or \Cref{eq:ns}).

\begin{figure}[H]
    \centering
    \includegraphics[width=0.6\linewidth]{example-image-a}
    \caption{Velocity contours at $t = \qty{10}{\second}$.}
    \label{fig:example}
\end{figure}

\section{References}
Citations use \texttt{biblatex}. As seen in \cite{versteeg2007}, the Finite Volume Method is standard for CFD.
